% Proved after four complete actual-file/source audits. Explicit failure of a uniform algebraic period shortcut; no geometric realization or main counting movement.
\subsection{An unbounded unitary scalar obstruction}

\begin{example}[Proved: no uniform scalar exponent from a unitary form]
\label{D:unitary-unbounded-scalar-obstruction}
Fix a prime $p$. Let $q$ be a power of $p$, and let
$m\ge3$ be odd with $m\mid q+1$.
Put $K=\mathbb F_{q^2}$ and let $\overline a=a^q$.
There is an injective homomorphism
\[
 \rho_m:(\mathbb Z/m)^2
       \lhook\joinrel\longrightarrow\operatorname{PSU}_m(q)
       \lhook\joinrel\longrightarrow\operatorname{PGL}_m(K)
                                                               \tag{D.UO.1}
\]
whose pullback of the scalar extension
\[
 1\longrightarrow K^\times\longrightarrow
 \operatorname{GL}_m(K)\longrightarrow
 \operatorname{PGL}_m(K)\longrightarrow1
                                                               \tag{D.UO.2}
\]
has class of exact order $m$ in
$H^2((\mathbb Z/m)^2,K^\times)$, with trivial
coefficient action. The same class is obtained by
restriction of the full projective-semilinear scalar
extension, with trivial field character.

For this fixed $p$, such odd $m$ are unbounded.
Consequently there is no fixed positive exponent
that kills the scalar obstruction of every map
into the natural projective unitary models as
rank and field vary. This remains so within the
nonabelian simple unitary ranges, by taking $m$
sufficiently large.
\end{example}

\begin{proof}
We use the standard unitary definitions and scalar
quotient of \citet[pp.~114--116, including Lemma~10.1
and its complete proof]{Taylor1992ClassicalB}.
In particular the norm-one subgroup of $K^\times$
is cyclic of order $q+1$. Choose $\zeta\in K^\times$
of exact order $m$; then $\zeta^{q+1}=1$.
Give $K^m$ the standard nondegenerate Hermitian form
\[
             H(v,w)=\sum_{j\in\mathbb Z/m}v_j\overline{w_j}.
                                                               \tag{D.UO.3}
\]
On its coordinate basis define the shift and diagonal
matrices
\[
 Xe_j=e_{j+1},\qquad Ye_j=\zeta^j e_j
                         \quad(j\in\mathbb Z/m).       \tag{D.UO.4}
\]
The definition of $Y$ is independent of the integer
representative of $j$. Clearly $X^m=Y^m=I$.
The matrix $X$ permutes an orthonormal basis, while
\[
 H(Yv,Yw)=\sum_j\zeta^j\overline{\zeta}^{\,j}
                              v_j\overline{w_j}
                         =H(v,w).
\]
Thus both matrices preserve the displayed Hermitian
form. Since $m$ is odd, their determinants are
\[
 \det X=(-1)^{m-1}=1,
 \qquad \det Y=\zeta^{m(m-1)/2}=1.                     \tag{D.UO.5}
\]
The last equality uses the integer $(m-1)/2$.
These formulas hold also in characteristic two.
Consequently $X,Y\in\operatorname{SU}_m(q)$.

A calculation on the coordinate basis gives
\[
       XY=\zeta^{-1}YX,
       \qquad XYX^{-1}Y^{-1}=\zeta^{-1}I.              \tag{D.UO.6}
\]
The scalar on the right belongs to the scalar
kernel of $\operatorname{SU}_m(q)\to
\operatorname{PSU}_m(q)$: it has norm one and its
$m$th power is one.
The projective images $\overline X,\overline Y$
therefore commute. If $X^aY^b$ is scalar for
$0\le a,b<m$, its permutation of the coordinate
lines forces $a=0$. The remaining diagonal matrix
has entries $1,\zeta^b,\ldots,\zeta^{(m-1)b}$;
it is scalar only if $\zeta^b=1$, hence $b=0$.
This proves that
$\rho_m(a,b)=\overline X^{\,a}\overline Y^{\,b}$
is the injection (D.UO.1), including its exact
projective kernel assertion.

Choose the set-theoretic lift
$s(a,b)=X^aY^b$ for $0\le a,b<m$.
There are no scalar factors from reducing exponents
modulo $m$, since $X^m=Y^m=I$.
Equation (D.UO.6) gives the scalar cocycle
\[
 \begin{aligned}
 s(a,b)s(a',b')
   &=c\bigl((a,b),(a',b')\bigr)s(a+a',b+b'),\\
 c\bigl((a,b),(a',b')\bigr)&=\zeta^{ba'}.
 \end{aligned}                                              \tag{D.UO.7}
\]
where the arguments of $s$ on the right are taken
modulo $m$. In particular $c^m=1$ pointwise, so
its class $\alpha_m$ is killed by $m$.

Conversely, every multiplicative coboundary for
the trivial action on the abelian group
$(\mathbb Z/m)^2$ is symmetric in its two arguments.
The alternating ratio of (D.UO.7), evaluated on
$g=(1,0)$ and $h=(0,1)$, is
\[
                   \frac{c(g,h)}{c(h,g)}=\zeta^{-1}.
                                                               \tag{D.UO.8}
\]
If $e\alpha_m=0$ for a positive integer $e$,
then $c^e$ is a coboundary and its alternating
ratio is one. Equation (D.UO.8) implies
$\zeta^e=1$, so $m\mid e$.
Together with the upper bound this proves that
$\alpha_m$ has exact order $m$ in the group with
coefficients $K^\times$ itself. No injectivity
from a smaller coefficient group is being assumed.
For the full semilinear extension the field
projection of $\rho_m$ is trivial, so restriction
recovers exactly (D.UO.2) and the same class.

For explicit unbounded choices at fixed $p$, if
$p=2$ take
\[
                       q=2^f,\qquad m=2^f+1
                       \quad(f\ge2).
                                                               \tag{D.UO.9}
\]
If $p$ is odd, take any positive even $f$ and put
\[
                       q=p^f,\qquad m=(p^f+1)/2.
                                                               \tag{D.UO.10}
\]
Here $p^f\equiv1\pmod8$, so $m$ is odd.
In both cases $m\mid q+1$, $p\nmid m$, and
$m$ tends to infinity. We may keep only values
$m\ge5$, so that the unitary groups lie in the
nonabelian simple ranges already specified in
Corollary~\ref{D:classical-defining-characteristic-crowns}.
The resulting exact orders $m$ rule out a uniform
scalar exponent.
\end{proof}

\begin{remark}[What this obstruction does and does not establish]
For the standard Hermitian form a scalar $uI$ has
multiplier $u^{q+1}$. That multiplier identity kills
the full scalar obstruction by $q+1$, and the
example shows that its consequence cannot be
replaced by an absolute exponent solely from the
existence of this natural unitary form.
The construction is an explicit algebraic test of
that shortcut, with no novelty claim.

No homomorphism from one fixed complete-local or
global geometric fundamental group is constructed
here. In particular there is no assertion that
$\alpha_m$ has a geometric Brauer image of order
$m$ on one fixed punctured surface, nor that an
unrestricted hereditary representation bound is
false. The result does not remove the unbounded-
period counting gate and is not a main reduction.
\end{remark}
